\section{上线策略：从小圈层到网络效应}

Elys 的早期传播来自小圈层，尤其是对 AI 新范式敏感的人群。这样的启动方式有好处：早期用户理解 context、分身和 proactive 的概念，也愿意容忍产品粗糙。坏处是小圈层不等于大众市场，破圈后隐私、误解、滥用和低质量互动都会放大。

\subsection{三阶段上线路线}

\noindent\begin{tabularx}{\textwidth}{p{0.22\textwidth}p{0.34\textwidth}X}
\toprule
阶段 & 目标 & 关键风险 \\
\midrule
种子圈层 & 验证范式和核心飞轮 & 用户过于懂 AI，不能代表大众。 \\
垂直社区 & 验证真人连接率和关系质量 & 社区边界、隐私和治理压力上升。 \\
大众市场 & 扩大网络效应和商业化 & AI 噪音、滥用、信任和监管风险。 \\
\bottomrule
\end{tabularx}

\begin{importantbox}{破圈不是单纯增长问题}
AI 社交破圈会改变风险结构。早期用户把分身当实验，大众用户可能把分身当本人。产品必须在破圈前建立权限、标记、申诉和治理机制。
\end{importantbox}

\subsection{本章小结}

Elys 这类产品不能只追求快速放量。AI 社交网络的上线策略必须和信任机制同步建设，否则网络效应会和风险一起增长。

